\lecture{24}{Wed 09 Nov 2022 10:31}{}

\begin{example}
	\begin{align*}
		\varphi: \mathbb{C} &\longrightarrow M_{2\times 2}(\mathbb{R}) \\
		a+bi &\longmapsto \varphi(a+bi) = 
	\left(\begin{array}{rr}a & b \\ -b & a\end{array}\right)
\end{align*}
This is a homomorphism of rings. We can verify that the map is injective.

So
\[
	\mathbb{C} \cong \varphi(\mathbb{C}) = \left\{
	\left(\begin{array}{rr}a & b \\ -b & a\end{array}\right):a,b \in \mathbb{R}
\right\} \subset M_{2\times 2}(\mathbb{R})
\]
\end{example}

\begin{theorem}[First Homomorphism Theorem]
	Let $\varphi: R \to R'$ be a surjective homomorphism of rings with $K = \operatorname{ker} \varphi$. Then $R / K' \cong R'$.
	In fact, 
	\begin{align*}
		\psi: R / K &\longrightarrow R' \\
		a + K &\longmapsto \psi(a + K) = \varphi(a)
	\end{align*}
	is an isomorphism.
\end{theorem}
\begin{proof}
	From the First Homomorphism Theorem of groups, we know that $\psi$ is an isomorphism of additive groups. It remains to show $\psi$ preserves multiplication. But
	\[
		\psi\left( (a+K)(b+K) \right) 
		=\psi\left( ab+K \right) 
		=\varphi(ab)
		=\varphi(a)\varphi(b)
		=\psi\left( a+K \right) \psi\left( b+K \right) 
	.\] 
\end{proof}
\begin{example}
	$\mathbb{Z}_p \cong \mathbb{Z} / p \mathbb{Z}$.
\end{example}

\begin{theorem}[Correspondence Theorem]
	Let $\varphi: R \to R'$ be surjective homomorphism of rings and $K = \operatorname{ker} \varphi$. Let $I' \triangleleft R'$. Let $I :=\varphi^{-1} (I') = \{a \in R: \varphi(a)\in I'\} $. Then $I \triangleleft R$ and $I' = \varphi(I)$. In fact this gives a 1-1 inclusion-preserving correspondence between ideals of $R$ containing $K$ and ideas of $R'$.
\end{theorem}
\begin{proof}
	We only need to check $I$, in addition to being a normal subgroup, is an ideal.
\end{proof}

\begin{theorem}[Second Homomorphism Theorem]
	Let $R$ be a ring, $A$ a subring and $I$ an ideal of $R$. Then $A + I$ is a subring of $R$ and $I$ is an ideal of $A + I$, and $I \cap A$ is an ideal of $A$. Also,
	\[
		\left( A + I \right)  / I \cong A / \left( A \cap I \right) 
	.\] 
\end{theorem}

\begin{theorem}[Third Homomorphism Theorem]
	If $K\triangleleft R$ and $I \triangleleft R$ and $I \supset K$, then
	\[
		R / K \cong \left( R / K \right) / (I / K)
	.\] 
\end{theorem}

\begin{example}
	\[
		K = p^2 \mathbb{Z} \triangleleft \mathbb{Z}, I = p\mathbb{Z} \triangleleft \mathbb{Z}, K \subset I 
	.\] 
	Hence
	\[
		\mathbb{Z} / p\mathbb{Z} \cong
		\left( \mathbb{Z} / p^2 \mathbb{Z} \right) / \left( p \mathbb{Z} / p^2 \mathbb{Z} \right) 
	.\] 
\end{example}
\begin{example}
	\begin{enumerate}
		\item Let $F$ be a field. Then $\{0\} $ and $F$ are the only ideals.
		\begin{proof}
			Let $I \neq \{0\}  $ be an ideal. Then there exists $a \in I$, $a \neq 0$. So there exists $a^{-1} \in F$. Now $1 = a^{-1} a \in I$ since $a \in I$ and $a^{-1}  \in F$. Let $r \in F$. Then $r = r 1 \in I$ since $1 \in I$. So $I = F$.
		\end{proof}
		\textbf{Consequence.} Let $\varphi: F \to R'$ be a homomorphism of rings, where $F$ is a field. Then $\operatorname{ker} \varphi \triangleleft F$. So $\varphi = 0$ ($\operatorname{ker} \varphi = F$ ) or $\varphi$ is injective ($\operatorname{ker} \varphi = \{0\} $ ). In the second case, $F \cong \varphi(F)$.
	\item Let $R $ be a commutative ring with $1$. Fix $a \in R$. Then $Ra = \{ra| r \in  R\} $. Notice $Ra \triangleleft R$ and $a = 1 a \in R a$. This is the unique smallest ideal containing $a$. This is called the \logseq{Principal Ideal} generated by $a$.
	\item Let $\mathbb{R}[t] = \{\text{polynomials with real coefficients}\} $ and 
		\begin{align*}
			\varphi: \mathbb{R}[t] &\longrightarrow \mathbb{C} \\
			f &\longmapsto \varphi(f) = f(i)
		.\end{align*}
		Then $\varphi$ is a surjective homomorphism of rings and $\operatorname{ker} \varphi = \mathbb{R}[t] (t^2+1)$.
		So by first homomorphism theorem,
		\[
			\mathbb{C} \cong\mathbb{R}[t] / \left( \mathbb{R}[t] (t^2 + 1) \right)
		.\]
		
	\end{enumerate}
\end{example}
